\chapter{Characters, Multiplicities, and the Center}
\label{ch:characters}

Keep the complex semisimple algebra and positive system of
Chapter~\ref{ch:highest-weights}. We use a nondegenerate invariant
form \(B\), positive on the real root space after identifying
\(\mathfrak h\) with its dual. In particular every squared norm
in a dominance argument is taken on that real space.
Let
\[
 \rho=\tfrac12\sum_{\alpha>0}\alpha,\quad
 \varepsilon(w)=\det(w),\quad
 A_\nu=\sum_{w\in W}\varepsilon(w)e^{w\nu},\quad
 D=e^\rho\prod_{\alpha>0}(1-e^{-\alpha}).
\]
The symbols \(e^\mu\) are formal basis elements, with
\(e^\mu e^\nu=e^{\mu+\nu}\), rather than a choice of a Lie-group
exponential.

\section{Formal characters and a finite truncation argument}

\begin{proposition}[Characters and their completion]
\label{res:H01}
For a finite-dimensional module set
\(\operatorname{ch}V=\sum_\mu(\dim V_\mu)e^\mu\in\mathbb Z[P]\).
Characters add in exact sequences, multiply in tensor products,
and are \(W\)-invariant. The characters of the \(L(\lambda)\),
\(\lambda\in P_+\), are linearly independent.
There is also a well-defined ring of series over the additive group
\(\mathfrak h^*\) whose support lies in a finite union of sets
\(\lambda-Q_+\). Integral-weight series form a subring; in the
larger ring
\[
 \operatorname{ch}M(\lambda)
   =e^\lambda\prod_{\alpha>0}(1-e^{-\alpha})^{-1}.
\]
\end{proposition}
\begin{proof}
Exactness holds weight by weight by
Lemma~\ref{supp:weight-submodules}, and
\((V\otimes U)_\eta=\bigoplus_{\mu+\nu=\eta}V_\mu\otimes U_\nu\).
These facts prove additivity and multiplication.
Weyl invariance is Proposition~\ref{res:G07}.
In a finite linear relation between simple characters, choose a
maximal highest weight among the terms. Its coefficient receives
no contribution from a different term, since all other weights
lie below that term's own highest weight. Its multiplicity in its
own character is one, so its coefficient is zero. Repeating proves
independence.

For the series ring, a coefficient of a product can receive only
finitely many contributions. Indeed, for fixed upper bounds
\(\lambda,\mu\), the decompositions
\(\beta+\gamma=\lambda+\mu-\eta\), \(\beta,\gamma\in Q_+\),
are finite: the simple-root coefficients are nonnegative integers
with bounded sum. A finite union of upper bounds makes no
difference to finiteness. The same observation proves associativity
by rearranging finite coefficient sums.
Finally PBW in \(U(\mathfrak n^-)\) makes a negative-root exponent
independent for every positive root. Summing each geometric series
gives the displayed formula. No Weyl action on this entire
completion is assumed; we use that action on finite expressions
when taking alternating sums.
\end{proof}

\begin{proposition}[The denominator, the Casimir, and truncated expansion]
\label{res:H02}
The denominator identity is \(D=A_\rho\).
The Casimir for \(B\) acts on any highest-weight module of weight
\(\lambda\) by
\[
 c_\lambda=(\lambda,\lambda+2\rho)
           =\|\lambda+\rho\|^2-\|\rho\|^2 .
\]
If a weight module \(V\), supported in \(\lambda-Q_+\), has finite
weight spaces and Casimir scalar \(c\), then, to any fixed height
bound, its character is an integral linear combination of
Verma characters \(M(\mu)\) with \(c_\mu=c\).
\end{proposition}
\begin{proof}
First \(s_i\) permutes the positive roots other than \(\alpha_i\),
and \(s_i\rho=\rho-\alpha_i\). Substitution in the product for \(D\)
gives \(s_iD=-D\).
In a finite alternating expression a coefficient on a reflection
hyperplane is zero. Every remaining orbit has a unique strictly
dominant representative \(\nu\), and its contribution is a
multiple of \(A_\nu\).
Every exponent of \(D\) is in \(\rho-Q_+\).
If such a \(\nu\) is strictly dominant and integral, then
\(\nu-\rho\) is dominant. If it were nonzero, its pairing with
\(\rho\) would be positive: \(\rho\) is a positive sum of simple
roots containing every simple root, and a nonzero dominant vector
pairs nonnegatively with all of them and positively with at least
one. On the other hand \(\nu-\rho\in-Q_+\), which pairs
nonpositively with \(\rho\), since
\((\rho,\alpha_i)>0\). This is impossible.
Thus only \(A_\rho\) occurs; its top coefficient is one.
This proves the denominator identity independently of the character
formula.

Choose dual bases \(h_a,h^a\) of \(\mathfrak h\), and root vectors
\(x_\alpha,y_\alpha\) with \(B(x_\alpha,y_\alpha)=1\).
Invariance gives
\([x_\alpha,y_\alpha]=t_\alpha\), where
\(B(t_\alpha,h)=\alpha(h)\).
The Casimir is
\[
 \Omega=\sum_a h_ah^a+
       \sum_{\alpha>0}(x_\alpha y_\alpha+y_\alpha x_\alpha)
 =\sum_a h_ah^a+2t_\rho+
       2\sum_{\alpha>0}y_\alpha x_\alpha.
\]
The dual tensor is invariant, so its multiplication commutes with
every \(x\in\mathfrak g\); this is also the calculation in
Theorem~\ref{res:C02}. Acting on a highest vector gives
\(\|\lambda\|^2+2(\rho,\lambda)\).
Since the module is generated by that vector and \(\Omega\) is
central, this is its action on the whole module.

For the expansion, the height of a weight \(\lambda-\beta\)
is \(\operatorname{ht}\beta\). If \(V\ne0\), take all weight
spaces of minimal height, and choose a basis of each.
There are finitely many of these spaces, each finite dimensional.
Every selected vector is highest, because raising a weight
would give a smaller height. The universal property gives a map
\[
 T=\bigoplus_j M(\mu_j)\longrightarrow V
\]
which is an isomorphism on the entire minimal-height layer.
The Casimir scalar on each summand is \(c_{\mu_j}=c\), as may
be checked on its nonzero image of the highest vector.
Both its kernel \(K\) and cokernel \(C\) therefore have Casimir
scalar \(c\) and have zero spaces on that layer. Exactness gives
\[
 \operatorname{ch}V=\operatorname{ch}T-\operatorname{ch}K+
                       \operatorname{ch}C.
\]
Repeat separately for \(K\) and \(C\). Their minimal heights have
strictly increased. To compute all coefficients through height
\(N\), stop after at most \(N+1\) such layers on each branch.
There are finitely many branches, and finitely many Verma modules
at each node. This proves the truncated assertion.
Taking the whole minimal layer is essential: removing only one
highest vector could leave the same highest weight in the cokernel.
The proof does not claim that an arbitrary module has a finite
untruncated Verma resolution.
\end{proof}

\begin{theorem}[Weyl character formula]
\label{res:H03}
For \(\lambda\in P_+\),
\[
 \operatorname{ch}L(\lambda)=
 \frac{A_{\lambda+\rho}}{A_\rho}
 =\frac{\sum_w\varepsilon(w)e^{w(\lambda+\rho)}}
        {e^\rho\prod_{\alpha>0}(1-e^{-\alpha})}.
\]
The fraction denotes the equality after multiplication by the
denominator, or its expansion in the completion of
Proposition~\ref{res:H01}.
\end{theorem}
\begin{proof}
The finite expression \(D\operatorname{ch}L(\lambda)\) is
alternating. Its exponents lie in \(\lambda+\rho-Q_+\), and the
coefficient of \(e^{\lambda+\rho}\) is one.
Apply the truncated expansion in Proposition~\ref{res:H02} at
arbitrarily large height bounds. Multiplication by
\(\prod_{\alpha>0}(1-e^{-\alpha})\) only lowers weights,
so any fixed coefficient is computed by a finite truncation.
Whenever \(e^{\mu+\rho}\) has a nonzero coefficient, some
Verma term contributing to it has \(c_\mu=c_\lambda\). Thus
\[
 \|\mu+\rho\|^2=\|\lambda+\rho\|^2.
\]
Choose the strictly dominant representative \(\nu\) of any orbit
in this alternating expression. It too is an exponent in the
support, so \(\delta=\lambda+\rho-\nu\in Q_+\).
If \(\delta\ne0\), strict dominance of both \(\lambda+\rho\)
and \(\nu\) gives
\[
 \|\lambda+\rho\|^2-\|\nu\|^2
     =(\delta,\lambda+\rho+\nu)>0,
\]
contrary to the Casimir equality. Therefore the sole possible
orbit is that of \(\lambda+\rho\).
Its top coefficient determines it to be \(A_{\lambda+\rho}\).
Singular orbits have zero coefficients by alternation.
This proves the formula without using Harish--Chandra's theorem.
\end{proof}

\begin{corollary}[Weyl dimension formula]
\label{res:H04}
For every dominant integral weight,
\[
 \dim L(\lambda)
  =\prod_{\alpha>0}\frac{(\lambda+\rho,\alpha)}{(\rho,\alpha)}
  =\prod_{\alpha>0}
       \frac{\langle\lambda+\rho,\alpha^\vee\rangle}
            {\langle\rho,\alpha^\vee\rangle}.
\]
The formula includes \(\lambda=0\), products of simple factors,
and the empty root system.
\end{corollary}
\begin{proof}
Evaluate the character at \(e^\mu=\exp(2t(\mu,\rho))\).
Invariance of the inner product and substitution \(w\mapsto w^{-1}\)
give
\[
 A_{\lambda+\rho}(2t\rho)
  =\sum_w\varepsilon(w)\exp(2t(w\rho,\lambda+\rho))
  =A_\rho(2t(\lambda+\rho)).
\]
By the denominator identity this is
\(\prod_{\alpha>0}
 (\exp(t(\alpha,\lambda+\rho))-\exp(-t(\alpha,\lambda+\rho)))\).
The denominator has the same expression with \(\lambda=0\).
Divide, and let \(t\to0\), using
\((e^{ta}-e^{-ta})/(e^{tb}-e^{-tb})\to a/b\) for \(b\ne0\).
Every denominator is positive by the choice of positive system.
The left side tends to the sum of the weight multiplicities.
Replacing roots by coroots cancels the same nonzero factor in
each quotient. All factors are independent on different simple
components, so the asserted conventions follow as well.
\end{proof}

\section{Weight multiplicities and tensor coefficients}

\begin{theorem}[Freudenthal's recursion]
\label{res:H05}
Let \(m_\lambda(\mu)=\dim L(\lambda)_\mu\), extended by zero
away from the weight set. Then
\[
 \bigl(\|\lambda+\rho\|^2-\|\mu+\rho\|^2\bigr)m_\lambda(\mu)
 =2\sum_{\alpha>0}\sum_{j\geq1}
        (\mu+j\alpha,\alpha)m_\lambda(\mu+j\alpha).
\]
Together with \(m_\lambda(\lambda)=1\), this is a terminating
recursion on the finite set
\((\lambda+Q)\cap\operatorname{conv}(W\lambda)\).
\end{theorem}
\begin{proof}
Set \(T_\mu^\alpha=\operatorname{tr}
(y_\alpha x_\alpha|_{L(\lambda)_\mu})\).
For maps between two different finite-dimensional spaces,
\(\operatorname{tr}(AB)=\operatorname{tr}(BA)\);
this follows by writing their two rectangular matrices.
Hence, using \([x_\alpha,y_\alpha]=t_\alpha\),
\[
 T_\mu^\alpha
   =\operatorname{tr}(x_\alpha y_\alpha|_{L(\lambda)_{\mu+\alpha}})
   =T_{\mu+\alpha}^\alpha+
                 (\mu+\alpha,\alpha)m_\lambda(\mu+\alpha).
\]
Iteration terminates because there are finitely many weights,
and gives
\(T_\mu^\alpha=\sum_{j\geq1}
(\mu+j\alpha,\alpha)m_\lambda(\mu+j\alpha)\).
Now take traces of the normal-ordered Casimir in
Proposition~\ref{res:H02} on the \(\mu\) weight space.
Its Cartan and linear terms have trace
\((\|\mu\|^2+2(\rho,\mu))m_\lambda(\mu)\).
Subtracting them from \(c_\lambda m_\lambda(\mu)\)
proves the identity.

By Proposition~\ref{res:G07} the weight set lies in the indicated
bounded polytope. For \(\mu\) in that polytope,
\(\|\mu\|^2\leq\|\lambda\|^2\), since a convex combination of
vectors of norm \(\|\lambda\|\) has no larger norm.
If also \(\mu\ne\lambda\) and \(\lambda-\mu\in Q_+\), then
\((\rho,\lambda-\mu)>0\). Consequently
\[
 \|\lambda+\rho\|^2-\|\mu+\rho\|^2
 =\|\lambda\|^2-\|\mu\|^2+2(\rho,\lambda-\mu)>0.
\]
All terms on the right have smaller height below \(\lambda\).
The lattice intersection with the polytope is finite. These
observations prove both that division is legitimate away from
the top weight and that the recursion terminates. Candidates
outside the cone or polytope have multiplicity zero.
\end{proof}

\begin{theorem}[Kostant's multiplicity formula]
\label{res:H06}
Let \(p(\beta)\) count the unordered decompositions of \(\beta\)
as a sum of positive roots, with \(p(0)=1\), and set \(p(\beta)=0\)
when \(\beta\notin Q_+\). Then
\[
 m_\lambda(\mu)=
   \sum_{w\in W}\varepsilon(w)
      p\bigl(w(\lambda+\rho)-(\mu+\rho)\bigr).
\]
\end{theorem}
\begin{proof}
Each root may occur with any nonnegative multiplicity, so in the
completed ring
\[
 \prod_{\alpha>0}(1-e^{-\alpha})^{-1}
     =\sum_{\beta\in Q_+}p(\beta)e^{-\beta}.
\]
Every coefficient is a finite sum by the height bound used in
Proposition~\ref{res:G02}.
Insert this expression into Theorem~\ref{res:H03}, cancel
\(e^\rho\), and compare the coefficient of \(e^\mu\).
The equality of exponents requires precisely
\(\beta=w(\lambda+\rho)-(\mu+\rho)\).
The zero convention covers weights outside the root-lattice
coset and outside the positive cone.
\end{proof}

\section{Polynomial invariants and central characters}

Polynomial functions on a vector space \(E\) form \(S(E^*)\).
The form \(B\) identifies \(S(\mathfrak g)\) with polynomial
functions on \(\mathfrak g\), compatibly with the adjoint action;
similarly on \(\mathfrak h\). When speaking of restriction we
always use these identifications. The group of inner automorphisms
is generated by \(\exp(t\operatorname{ad}x)\), \(x\in\mathfrak g\).
For root vectors these exponentials are finite on \(\mathfrak g\);
their rank-one products represent the Weyl reflections.

\begin{theorem}[Chevalley restriction]
\label{res:H07}
Restriction is an isomorphism
\[
 S(\mathfrak g^*)^{\mathfrak g}\longrightarrow
 S(\mathfrak h^*)^W .
\]
The source is generated as an algebra by trace polynomials
\(x\mapsto\operatorname{tr}_{L(\lambda)}(x^m)\),
\(\lambda\in P_+\), \(m\geq1\), and constants.
\end{theorem}
\begin{proof}
Infinitesimal invariance is equivalent to invariance under all
inner exponentials. If every infinitesimal action kills a
polynomial, its value along \(\exp(t\operatorname{ad}x)y\) has
zero derivative for every \(t\), and hence is constant.
Conversely differentiate such invariance at \(t=0\).
The root directions suffice to test infinitesimal invariance,
because their operators generate the Lie algebra action.
The representatives of the simple Weyl reflections are products
of these exponentials, so restriction lands in the stated
invariants. If restriction is zero, invariance makes the polynomial
vanish on all conjugates of \(\mathfrak h\).
They include the dense regular semisimple set from
Proposition~\ref{res:D06}; a polynomial vanishing on a dense set
is zero. This proves injectivity.

Here is a direct surjectivity proof, including the invariant
polynomials which must be constructed. The lattice \(P\) spans
\(\mathfrak h^*\), and pure powers of lattice elements span its
symmetric powers. For example polarization expresses a product
\(\eta_1\cdots\eta_m\) as
\[
 \frac1{m!}\sum_{J\subseteq\{1,\ldots,m\}}
       (-1)^{m-|J|}\left(\sum_{j\in J}\eta_j\right)^m.
\]
Averaging over \(W\), and then moving a lattice element to the
dominant chamber, shows that the invariant polynomials are spanned
by orbit sums
\(q_{\lambda,m}=\sum_{\nu\in W\lambda}\nu^m\),
\(\lambda\in P_+\). Constants can be treated separately.

For such a \(\lambda\), take \(L(\lambda)\), whose existence has
already been proved. The trace polynomial
\(T_{\lambda,m}(x)=\operatorname{tr}(\pi_\lambda(x)^m)\)
is infinitesimally invariant: its derivative in the direction
\([y,x]\) is
\(m\operatorname{tr}(\pi_\lambda(x)^{m-1}
[\pi_\lambda(y),\pi_\lambda(x)])=0\) by trace cyclicity.
The preceding paragraph gives its inner invariance as well.
Its restriction is
\[
 T_{\lambda,m}|_{\mathfrak h}
   =q_{\lambda,m}
     +\sum_{\substack{\nu\in P_+\\\nu<\lambda}}
                  m_\lambda(\nu)q_{\nu,m}.
\]
Only weights of \(L(\lambda)\) occur. There are finitely many
dominant weights below any fixed \(\lambda\), by the
\(\rho^\vee\) bound in Theorem~\ref{res:G05}.
Induction on that finite lower set therefore lifts every
\(q_{\lambda,m}\) as a linear combination of trace polynomials.
This proves surjectivity and the generation assertion.
No conclusion about surjectivity has been inferred just from
root-direction differentiation.
\end{proof}

\begin{theorem}[Harish--Chandra isomorphism and central characters]
\label{res:H08}
Let \(Z=Z(U(\mathfrak g))\). In the triangular PBW ordering,
project onto the monomials entirely in \(\mathfrak h\):
\(p:U(\mathfrak g)\to U(\mathfrak h)=S(\mathfrak h)\).
For \(z\in Z\), define
\[
 \gamma(z)(\eta)=p(z)(\eta-\rho),\qquad \eta\in\mathfrak h^*.
\]
Then \(\gamma:Z\longrightarrow S(\mathfrak h)^W\) is an
algebra isomorphism. The central character on \(M(\lambda)\)
and its nonzero highest-weight quotients is
\(\chi_\lambda(z)=\gamma(z)(\lambda+\rho)\), and
\[
 \chi_\lambda=\chi_\mu
    \quad\Longleftrightarrow\quad
 \mu+\rho\in W(\lambda+\rho).
\]
Every complex algebra character of \(Z\) is obtained in this way.
\end{theorem}
\begin{proof}
A central \(z\) has Cartan weight zero. In its PBW expansion,
a term with a positive-root factor on the right kills a highest
vector. A term with negative-root factors and none positive
cannot have weight zero unless those negative factors are absent.
Thus \(zv_\lambda=p(z)(\lambda)v_\lambda\).
Centrality extends this scalar action to the generated module.
It follows for all \(\lambda\) that
\(p(zz')(\lambda)=p(z)(\lambda)p(z')(\lambda)\).
Two polynomials equal at all \(\lambda\) are equal, proving
multiplicativity on \(Z\).

We check the shifted invariance without presuming the theorem.
If \(\lambda(h_i)=m\in\mathbb Z_{\geq0}\), the vector
\(f_i^{m+1}v_\lambda\) in \(M(\lambda)\) is nonzero by PBW.
The rank-one relation makes it killed by \(e_i\), and it is
killed by \(e_j\), \(j\ne i\), because \([e_j,f_i]=0\).
It has weight \(s_i(\lambda+\rho)-\rho\).
There is therefore a nonzero highest-weight map from the Verma
module of that weight into \(M(\lambda)\).
The center has the same scalar on the two modules, so
\[
 p(z)(\lambda)
    =p(z)\bigl(s_i(\lambda+\rho)-\rho\bigr).
\]
For each nonnegative integer \(m\) this is a polynomial identity
on the entire affine hyperplane \(\lambda(h_i)=m\);
the other coordinates are arbitrary. A polynomial in
\(\lambda(h_i)\), with coefficients in the remaining coordinates,
cannot have infinitely many distinct roots unless it is zero.
The identity thus holds everywhere. After translation by
\(\rho\) it says that \(\gamma(z)\) is \(s_i\)-invariant,
and hence \(W\)-invariant.

Equip \(U(\mathfrak g)\) with PBW degree.
The leading symbol of a central element is in \(S(\mathfrak g)^{\mathfrak g}\).
The leading symbol of \(p(z)\), if not zero, is its restriction
to \(\mathfrak h\). By Chevalley restriction that restriction
cannot be zero for a nonzero invariant symbol. Translation
by \(-\rho\) does not change the leading symbol. It follows that
\(\gamma(z)\ne0\) for \(z\ne0\), proving injectivity.

For surjectivity, take a homogeneous invariant polynomial
\(q\in S(\mathfrak h)^W\) of degree \(d\).
By Theorem~\ref{res:H07} and the form \(B\), lift it to a
homogeneous \(F\in S(\mathfrak g)^{\mathfrak g}\).
The symmetrization map
\[
 x_1\cdots x_d\longmapsto
       \frac1{d!}\sum_{\sigma\in S_d}x_{\sigma(1)}\cdots x_{\sigma(d)}
\]
is a filtered \(\mathfrak g\)-module isomorphism to
\(U(\mathfrak g)\) by PBW.
It carries \(F\) to a central element \(z\) whose leading
symbol restricts to \(q\). Thus \(\gamma(z)-q\) is an invariant
of smaller degree. Induction on degree corrects this difference;
degree zero is immediate. This proves the isomorphism.

Finally finite-group invariants separate distinct orbits.
On two disjoint finite orbits choose, by polynomial interpolation,
a polynomial equal to one on the first and zero on the second.
One construction makes a delta polynomial at each point by
multiplying linear functions separating it from all other points,
then sums those delta polynomials. Averaging over \(W\) keeps
the two prescribed constant values and makes it invariant.
This proves the equivalence of central characters.
To see that no characters are missing, the polynomial ring
\(A=S(\mathfrak h)\) is integral and finite over \(A^W\):
each of its finitely many linear generators satisfies the monic
polynomial \(\prod_{w\in W}(t-w\ell)\), and adjoining finitely
many integral generators is finite.
Lying-over lifts a maximal ideal of \(A^W\) to one of \(A\).
The weak Nullstellensatz identifies that ideal of \(A\) with
evaluation at some \(\eta\in\mathfrak h^*\).
These elementary integral-extension facts are among the
prerequisites in \cite[Chapter 4]{p103}.
Take \(\lambda=\eta-\rho\).
\end{proof}

\begin{theorem}[Steinberg's tensor multiplicities and small-rank checks]
\label{res:H09}
If \(\lambda,\mu,\nu\in P_+\), the multiplicity of \(L(\nu)\)
in \(L(\lambda)\otimes L(\mu)\) is
\[
 c_{\lambda\mu}^{\nu}
 =\sum_{u,w\in W}\varepsilon(u)\varepsilon(w)
   p\bigl(u(\lambda+\rho)+w(\mu+\rho)-(\nu+2\rho)\bigr).
\]
For type \(A_1\) this gives
\(L(a)\otimes L(b)=\bigoplus_{j=0}^{\min(a,b)}L(a+b-2j)\).
For \(A_2\),
\[
 \dim L(a\omega_1+b\omega_2)
    =\frac{(a+1)(b+1)(a+b+2)}2 ,
 \qquad p(s\alpha_1+t\alpha_2)=
      \min(s,t)+1\quad(s,t\geq0).
\]
\end{theorem}
\begin{proof}
By complete reducibility and the character formula,
\[
 \operatorname{ch}L(\lambda) A_{\mu+\rho}
       =\sum_{\xi\in P_+}c_{\lambda\mu}^{\xi}A_{\xi+\rho}.
\]
For dominant \(\nu\), the coefficient of \(e^{\nu+\rho}\)
on the right is just \(c_{\lambda\mu}^\nu\):
two strictly dominant weights in the same Weyl orbit agree.
The coefficient on the left is
\(\sum_w\varepsilon(w)m_\lambda(\nu+\rho-w(\mu+\rho))\).
Insert Kostant's formula and simplify its argument to obtain
the stated double sum.

For \(A_1\), the rank-one character is
\(e^a+e^{a-2}+\cdots+e^{-a}\).
The coefficient counts in the product, followed by successive
subtraction of its top rank-one character, give the indicated
list, each once. This also checks the double sum since
\(p(j\alpha)=1\) for \(j\geq0\) and zero otherwise.
For \(A_2\) the positive roots are
\(\alpha_1,\alpha_2,\alpha_1+\alpha_2\).
The dimension product has factors
\(a+1,b+1,(a+b+2)/2\).
A decomposition of \(s\alpha_1+t\alpha_2\) uses
\(j\) copies of \(\alpha_1+\alpha_2\), with
\(0\leq j\leq\min(s,t)\), and then forced numbers of the
two simple roots. This proves the partition formula.
As an internal multiplicity check, for the adjoint highest
weight \(\lambda=\alpha_1+\alpha_2\), Kostant's formula at
\(\mu=0\) has the identity term \(p(\alpha_1+\alpha_2)=2\);
every other Weyl term lies outside \(Q_+\).
It gives the two-dimensional zero weight space \(\mathfrak h\),
consistent with the eight-dimensional decomposition in
Example~\ref{res:G09}.
\end{proof}
